\section{未来挑战}

\subsection{可解释与审计的灵活组合}

随着模型越来越强，单纯的 replay evaluation 不再能保证安全，必须引入 mechanistic interpretability 和 audit trails。

\begin{itemize}
    \item 记录每个 action 触发的 Q gradient 信息，便于 reverse engineer。
    \item 将 audit log 与 governance checklist 搭配，形成 "explain -> sign off" 循环。
    \item 在 policy update 时触发 automated sanity checks 并留存 diff。
\end{itemize}

\begin{warningbox}{忽略解释性会导致合规缺失}
若没有办法解释为什么 policy 选择某个动作，治理团队会将其标记为“不透明”，从而延迟部署。解释性与 compliance 是两个互为支撑的维度。
\end{warningbox}

\subsection{适应法规与政策变更}

讲者强调政策变化对 offline RL 项目影响巨大：数据共享限制、logging 要求、甚至禁止某些 reward shaping 方式。需要实时跟踪法规并将文档化变更归档。

\subsection{本章小结}

未来的 offline RL 项目必须把 interpretability 和 governance 绑在一起，并保持政策文档与技术文档同步，才能在监管压力下稳健前行。

